\chapter{Complexity — Random Walk Mixing}
%%% generated by the blueprint pipeline from TCSlib.Complexity.Randomized.RandomWalkMixing
\section{Overview}

Declarations of the Lean module \texttt{TCSlib.Complexity.Randomized.RandomWalkMixing}.

\section{Declarations}

\begin{definition}[Maximum excess over uniformity]
\lean{Randomized.uniformExcess}
\statementsource{levin-peres-wilmer-ch04-mixing}{pdf-70350d84fb7f-p002-b005}
\label{Randomized.uniformExcess}
\leanok
Let $V$ be a finite nonempty set and let $p:V\to\mathbb{Q}$. The uniform excess of $p$
is the maximum, over $i\in V$, of $p(i)-1/|V|$.
\end{definition}

\begin{theorem}[Uniform excess bound]
\lean{Randomized.le_uniformExcess}
\label{Randomized.le_uniformExcess}
\leanok
\uses{Randomized.uniformExcess}
\difficulty{1}
Let $V$ be a finite nonempty set, let $p:V\to\mathbb{Q}$, and let $i\in V$. Then
$p(i)-1/|V|$ is at most the maximum of $p(v)-1/|V|$ over $v\in V$.
\end{theorem}

\begin{theorem}[Centered probabilities sum to zero]
\lean{Randomized.sum_sub_inv_card}
\label{Randomized.sum_sub_inv_card}
\leanok
\difficulty{3}
Let $V$ be a finite nonempty set and let $p:V\to\mathbb{Q}$ satisfy $\sum_{v\in
V}p(v)=1$. Then $\sum_{v\in V}\bigl(p(v)-1/|V|\bigr)=0$.
\end{theorem}

\begin{theorem}[Nonnegativity of uniform excess]
\lean{Randomized.uniformExcess_nonneg}
\label{Randomized.uniformExcess_nonneg}
\leanok
\uses{Randomized.le_uniformExcess, Randomized.sum_sub_inv_card, Randomized.uniformExcess}
\difficulty{4}
Let $V$ be a finite nonempty set and let $p:V\to\mathbb{Q}$ have sum $1$. The maximum of
$p(v)-1/|V|$ over $v\in V$ is nonnegative.
\end{theorem}

\begin{theorem}[Lower bound from uniform excess]
\lean{Randomized.inv_card_sub_le}
\label{Randomized.inv_card_sub_le}
\leanok
\uses{Randomized.le_uniformExcess, Randomized.sum_sub_inv_card, Randomized.uniformExcess, Randomized.uniformExcess_nonneg}
\difficulty{5}
Let $V$ be a finite nonempty set, and let $p:V\to\mathbb{Q}$ satisfy $\sum_{u\in
V}p(u)=1$. Define the uniform excess to be $\max_{u\in V}\bigl(p(u)-1/|V|\bigr)$. Then,
for every $v\in V$, $1/|V|-|V|\max_{u\in V}\bigl(p(u)-1/|V|\bigr)\le p(v)$.
\end{theorem}

\begin{theorem}[Uniform excess under a minorized transition]
\lean{Randomized.uniformExcess_le_of_minorant}
\label{Randomized.uniformExcess_le_of_minorant}
\leanok
\uses{Randomized.le_uniformExcess, Randomized.sum_sub_inv_card, Randomized.uniformExcess}
\difficulty{6}
Let $V$ be a finite nonempty set, let $P:V\times V\to\mathbb{Q}$ have each column sum
equal to $1$, and suppose $P(u,v)\ge\delta$ for every $u,v\in V$, where
$\delta\in\mathbb{Q}$. Let $p:V\to\mathbb{Q}$ have sum $1$. For a function
$q:V\to\mathbb{Q}$, define its uniform excess to be $\max_{v\in V}(q(v)-1/|V|)$. Then
the uniform excess of $v\mapsto\sum_{u\in V}p(u)P(u,v)$ is at most $(1-|V|\delta)$ times
the uniform excess of $p$.
\end{theorem}

\begin{definition}[Transition probability]
\lean{Randomized.RegularMultigraph.transProb}
\statementsource{Vad12}{
  pdf-bf2ce5422088-p039-b002,
  pdf-bf2ce5422088-p039-b003
}
\statementsource{arora-barak-ch07-randomized-computation}{pdf-26d95080dbde-p021-b002}
\label{Randomized.RegularMultigraph.transProb}
\leanok
\uses{Randomized.RegularMultigraph.endpoint, Randomized.idxProb}
Let $G$ be a regular multigraph of degree $D$ on a vertex type $V$. For vertices $u,v\in
V$ and an integer $k\ge 0$, the transition probability is the index-sampling probability
that the endpoint reached from $u$ by index choice $w$ equals $v$.
\end{definition}

\begin{theorem}[Zero-step transition probability]
\lean{Randomized.RegularMultigraph.transProb_zero}
\label{Randomized.RegularMultigraph.transProb_zero}
\leanok
\uses{Randomized.RegularMultigraph.transProb, Randomized.idxProb_zero}
\difficulty{3}
Let $G$ be a regular multigraph of degree $D$, and let $u,v$ be vertices of $G$. The
probability that a zero-step walk from $u$ ends at $v$ is $1$ if $u=v$ and $0$
otherwise.
\end{theorem}

\begin{theorem}[Recurrence for transition probabilities]
\lean{Randomized.RegularMultigraph.transProb_succ}
\label{Randomized.RegularMultigraph.transProb_succ}
\leanok
\uses{Randomized.RegularMultigraph.avg, Randomized.RegularMultigraph.transProb, Randomized.idxProb_succ}
\difficulty{3}
Let $G$ be a regular multigraph of positive degree $D$, let $k$ be a nonnegative
integer, and let $u,v$ be vertices of $G$. The probability that a walk of $k+1$ steps
from $u$ ends at $v$ equals the average, over the $D$ neighbor indices at $u$, of the
probability that a walk of $k$ steps from the resulting neighbor ends at $v$.
\end{theorem}

\begin{theorem}[One-step transition probability]
\lean{Randomized.RegularMultigraph.transProb_one}
\label{Randomized.RegularMultigraph.transProb_one}
\leanok
\uses{Randomized.RegularMultigraph.avg, Randomized.RegularMultigraph.mult, Randomized.RegularMultigraph.transProb, Randomized.RegularMultigraph.transProb_succ, Randomized.RegularMultigraph.transProb_zero}
\difficulty{3}
Let $G$ be a regular multigraph of positive degree $D$, and let $u,v$ be vertices of
$G$. The probability that a one-step walk from $u$ ends at $v$ is the multiplicity of
the edge from $u$ to $v$, divided by $D$.
\end{theorem}

\begin{theorem}[Nonnegativity of walk transition probabilities]
\lean{Randomized.RegularMultigraph.transProb_nonneg}
\label{Randomized.RegularMultigraph.transProb_nonneg}
\leanok
\uses{Randomized.RegularMultigraph.transProb, Randomized.idxProb}
\difficulty{2}
Let $G$ be a regular multigraph of degree $D$, and let $k$ be a nonnegative integer and
$u,v$ vertices of $G$. The probability that a walk starting at $u$, with its successive
neighbor indices chosen uniformly from the $D$ possibilities, ends at $v$ after $k$
steps is nonnegative.
\end{theorem}

\begin{theorem}[Lower bound for a prescribed walk]
\lean{Randomized.RegularMultigraph.one_div_pow_le_transProb}
\label{Randomized.RegularMultigraph.one_div_pow_le_transProb}
\leanok
\uses{Randomized.RegularMultigraph.endpoint, Randomized.RegularMultigraph.transProb, Randomized.idxProb}
\difficulty{4}
Let $G$ be a regular multigraph of positive degree $D$, and let $u,v$ be vertices. If a
sequence $w$ of neighbor indices in $\mathrm{Fin}(D)$ takes the walk from $u$ to $v$,
then the probability that a uniformly chosen sequence of length $|w|$ takes the walk
from $u$ to $v$ is at least $D^{-|w|}$.
\end{theorem}

\begin{definition}[Number of visits along a walk]
\lean{Randomized.RegularMultigraph.visits}
\statementsource{Vad12}{pdf-bf2ce5422088-p038-b006}
\statementsource{arora-barak-ch07-randomized-computation}{pdf-26d95080dbde-p016-b002}
\label{Randomized.RegularMultigraph.visits}
\leanok
Let $G$ be a regular multigraph of degree $D$ on a vertex type $V$. For vertices $v,u\in
V$ and a finite list of edge indices, the number of visits to $v$ is defined
recursively: the empty list gives zero, and for a list beginning with index $i$, count
one if the neighbor of $u$ at index $i$ is $v$, then add the visits to $v$ along the
remaining list starting from that neighbor.
\end{definition}

\begin{definition}[Expected number of visits]
\lean{Randomized.RegularMultigraph.expVisits}
\statementsource{arora-barak-ch07-randomized-computation}{pdf-26d95080dbde-p016-b002}
\label{Randomized.RegularMultigraph.expVisits}
\leanok
\uses{Randomized.RegularMultigraph.visits}
Let $G$ be a regular multigraph of degree $D$ on a vertex type $V$. For vertices $v,u\in
V$ and an integer $N\ge 0$, the expected number of visits to $v$ in an $N$-step walk
from $u$ is the sum of the visit counts over all sequences of $N$ indices in
$\mathrm{Fin}(D)$, divided by $D^N$.
\end{definition}

\begin{theorem}[Expected Visit Recursion]
\lean{Randomized.RegularMultigraph.expVisits_succ}
\label{Randomized.RegularMultigraph.expVisits_succ}
\leanok
\uses{Randomized.RegularMultigraph.expVisits, Randomized.RegularMultigraph.visits}
\difficulty{5}
Let $G$ be a regular multigraph of degree $D>0$, and let $u$ and $v$ be vertices. For
every natural number $N$, the expected number of visits to $v$ in $N+1$ steps starting
at $u$ equals the average, over the $D$ possible first edge indices, of the indicator
that the first step reaches $v$ plus the expected number of visits to $v$ in the
remaining $N$ steps starting at the resulting neighbor.
\end{theorem}

\begin{theorem}[Expected visits and transition probabilities]
\lean{Randomized.RegularMultigraph.expVisits_eq}
\label{Randomized.RegularMultigraph.expVisits_eq}
\leanok
\uses{Randomized.RegularMultigraph.avg, Randomized.RegularMultigraph.expVisits, Randomized.RegularMultigraph.expVisits_succ, Randomized.RegularMultigraph.transProb, Randomized.RegularMultigraph.transProb_succ, Randomized.RegularMultigraph.transProb_zero, Randomized.RegularMultigraph.visits}
\difficulty{5}
Let $G$ be a regular multigraph of degree $D>0$, and let $u$ and $v$ be vertices. For
every $N\in\mathbb{N}$, the expected number of visits to $v$ along a uniformly chosen
sequence of $N$ neighbor indices starting at $u$ equals the sum, over $j=0,\ldots,N-1$,
of the probabilities that the walk is at $v$ after $j+1$ steps.
\end{theorem}

\begin{theorem}[Total transition probability from a vertex]
\lean{Randomized.RegularMultigraph.sum_transProb_right}
\label{Randomized.RegularMultigraph.sum_transProb_right}
\leanok
\uses{Randomized.RegularMultigraph.avg, Randomized.RegularMultigraph.transProb, Randomized.RegularMultigraph.transProb_succ, Randomized.RegularMultigraph.transProb_zero}
\difficulty{4}
Let $G$ be a regular multigraph on a finite vertex set, with positive degree $D$. For
every nonnegative integer $k$ and every vertex $u$, the sum over all vertices $v$ of the
probability that a uniformly chosen sequence of $k$ neighbor indices, followed from $u$,
ends at $v$ is $1$.
\end{theorem}

\begin{theorem}[Total transition probability into a vertex]
\lean{Randomized.RegularMultigraph.sum_transProb_left}
\label{Randomized.RegularMultigraph.sum_transProb_left}
\leanok
\uses{Randomized.RegularMultigraph.sum_avg, Randomized.RegularMultigraph.transProb, Randomized.RegularMultigraph.transProb_succ, Randomized.RegularMultigraph.transProb_zero}
\difficulty{4}
Let $G$ be a finite regular multigraph of positive degree $D$, and let $k$ be a
nonnegative integer. For any vertex $v$, the sum over all starting vertices of the
probability that a uniformly chosen sequence of $k$ neighbor indices ends at $v$ is $1$.
\end{theorem}

\begin{theorem}[Transition probabilities summed over intermediate vertices]
\lean{Randomized.RegularMultigraph.transProb_succ_eq_sum}
\label{Randomized.RegularMultigraph.transProb_succ_eq_sum}
\leanok
\uses{Randomized.RegularMultigraph.avg, Randomized.RegularMultigraph.sum_nbr_eq, Randomized.RegularMultigraph.transProb, Randomized.RegularMultigraph.transProb_one, Randomized.RegularMultigraph.transProb_succ}
\difficulty{3}
Let $G$ be a finite regular multigraph of positive degree $D$, and let $k$ be a
nonnegative integer and $u,v$ vertices of $G$. The probability that a walk of $k+1$
steps from $u$ ends at $v$ equals the sum over all vertices $x$ of the product of the
one-step transition probability from $u$ to $x$ and the $k$-step transition probability
from $x$ to $v$.
\end{theorem}

\begin{definition}[Distribution after a walk]
\lean{Randomized.RegularMultigraph.walkDist}
\statementsource{Vad12}{pdf-bf2ce5422088-p039-b003}
\statementsource{arora-barak-ch07-randomized-computation}{pdf-26d95080dbde-p021-b002}
\label{Randomized.RegularMultigraph.walkDist}
\leanok
\uses{Randomized.RegularMultigraph.transProb}
Let $G$ be a regular multigraph of degree $D$ on a finite vertex type $V$, and let
$p:V\to\mathbb{Q}$. For an integer $k\ge 0$ and vertex $v\in V$, the walk distribution
at $v$ is $\sum_{u\in V}p(u)\,T_k(u,v)$, where $T_k(u,v)$ is the transition probability
from $u$ to $v$.
\end{definition}

\begin{theorem}[Zero-Step Walk Distribution]
\lean{Randomized.RegularMultigraph.walkDist_zero}
\label{Randomized.RegularMultigraph.walkDist_zero}
\leanok
\uses{Randomized.RegularMultigraph.transProb_zero, Randomized.RegularMultigraph.walkDist}
\difficulty{3}
Let $G$ be a regular multigraph on a finite vertex set, and let $p$ assign a rational
weight to each vertex. Walking for zero steps from $p$ leaves the distribution
unchanged.
\end{theorem}

\begin{theorem}[Successive Walk Distribution]
\lean{Randomized.RegularMultigraph.walkDist_succ}
\label{Randomized.RegularMultigraph.walkDist_succ}
\leanok
\uses{Randomized.RegularMultigraph.transProb_succ_eq_sum, Randomized.RegularMultigraph.walkDist}
\difficulty{4}
Let $G$ be a regular multigraph of positive degree $D$ on a finite vertex set, let $p$
assign a rational weight to each vertex, and let $k$ be a nonnegative integer. The
distribution obtained by walking for $k+1$ steps from $p$ equals the distribution
obtained by first walking for one step and then for $k$ steps.
\end{theorem}

\begin{theorem}[Composition of Walk Distributions]
\lean{Randomized.RegularMultigraph.walkDist_add}
\label{Randomized.RegularMultigraph.walkDist_add}
\leanok
\uses{Randomized.RegularMultigraph.walkDist, Randomized.RegularMultigraph.walkDist_succ, Randomized.RegularMultigraph.walkDist_zero}
\difficulty{4}
Let $G$ be a regular multigraph of positive degree $D$ on a finite vertex set, let $p$
assign a rational weight to each vertex, and let $a,b$ be nonnegative integers. The
distribution obtained by walking for $a+b$ steps from $p$ equals the distribution
obtained by first walking for $a$ steps and then walking for $b$ steps.
\end{theorem}

\begin{theorem}[Conservation of walk distribution mass]
\lean{Randomized.RegularMultigraph.sum_walkDist}
\label{Randomized.RegularMultigraph.sum_walkDist}
\leanok
\uses{Randomized.RegularMultigraph.sum_transProb_right, Randomized.RegularMultigraph.walkDist}
\difficulty{3}
Let $G$ be a finite regular multigraph of positive degree $D$, let $p$ be any
rational-valued function on its vertices, and let $k$ be a nonnegative integer. The sum
over destination vertices of the walk distribution obtained after $k$ steps from $p$
equals the sum of $p$ over all vertices.
\end{theorem}

\begin{definition}[Stationary distribution]
\lean{Randomized.RegularMultigraph.IsStationary}
\statementsource{Vad12}{pdf-bf2ce5422088-p039-b003}
\statementsource{arora-barak-ch07-randomized-computation}{pdf-26d95080dbde-p019-b001}
\label{Randomized.RegularMultigraph.IsStationary}
\leanok
\uses{Randomized.RegularMultigraph.walkDist}
Let $G$ be a regular multigraph of degree $D$ on a finite vertex type $V$, and let
$p:V\to\mathbb{Q}$. The function $p$ is stationary if its distribution after one step of
the walk equals $p$.
\end{definition}

\begin{theorem}[Uniform Stationarity]
\lean{Randomized.RegularMultigraph.uniform_isStationary}
\statementsource{Gil98}{pdf-d5b902970e61-p005-b001}
\label{Randomized.RegularMultigraph.uniform_isStationary}
\leanok
\uses{Randomized.RegularMultigraph.IsStationary, Randomized.RegularMultigraph.transProb_one, Randomized.RegularMultigraph.uniform_stationary, Randomized.RegularMultigraph.walkDist}
\difficulty{3}
Let $G$ be a finite regular multigraph of degree $D>0$. The uniform distribution on its
vertices is stationary for the random walk that at each step chooses one of the $D$
indexed neighbors uniformly.
\end{theorem}

\begin{theorem}[A lower bound on transition probabilities]
\lean{Randomized.RegularMultigraph.one_div_pow_le_transProb_two_mul}
\label{Randomized.RegularMultigraph.one_div_pow_le_transProb_two_mul}
\leanok
\uses{Randomized.RegularMultigraph.Connected, Randomized.RegularMultigraph.IsBipartite, Randomized.RegularMultigraph.exists_endpoint_eq_two_mul, Randomized.RegularMultigraph.one_div_pow_le_transProb, Randomized.RegularMultigraph.transProb}
\difficulty{4}
Let $G$ be a finite regular multigraph of positive degree $D$ that is connected and not
bipartite. For any vertices $u,v$, the probability that a uniformly chosen sequence of
$2|V|$ neighbor indices takes the walk from $u$ to $v$ is at least $D^{-2|V|}$, where
$|V|$ is the number of vertices.
\end{theorem}

\begin{theorem}[Uniform excess walk dist le self]
\lean{Randomized.RegularMultigraph.uniformExcess_walkDist_le_self}
\statementsource{levin-peres-wilmer-ch04-mixing}{pdf-70350d84fb7f-p011-b004}
\label{Randomized.RegularMultigraph.uniformExcess_walkDist_le_self}
\leanok
\uses{Randomized.RegularMultigraph.sum_transProb_left, Randomized.RegularMultigraph.transProb, Randomized.RegularMultigraph.transProb_nonneg, Randomized.RegularMultigraph.walkDist, Randomized.uniformExcess, Randomized.uniformExcess_le_of_minorant}
\difficulty{3}
% dataset-pipeline: stub -- prose pending build_informal_statements.py
Statement pending regeneration.
\end{theorem}

\begin{theorem}[Iterated contraction of uniform excess]
\lean{Randomized.RegularMultigraph.uniformExcess_walkDist_mul_le}
\label{Randomized.RegularMultigraph.uniformExcess_walkDist_mul_le}
\leanok
\uses{Randomized.RegularMultigraph.sum_walkDist, Randomized.RegularMultigraph.walkDist, Randomized.RegularMultigraph.walkDist_add, Randomized.RegularMultigraph.walkDist_zero, Randomized.uniformExcess}
\difficulty{4}
Let $G$ be a finite regular multigraph of positive degree $D$ with a nonempty vertex
set. Suppose $c$ is a nonnegative rational number such that, for every rational-valued
function $p$ on the vertices whose values sum to $1$, applying the $K$-step walk
distribution to $p$ gives uniform excess at most $c$ times the uniform excess of $p$.
Then for every nonnegative integer $i$ and every such function $p$, the distribution
after $iK$ steps has uniform excess at most $c^i$ times the uniform excess of $p$.
\end{theorem}

\begin{theorem}[Auxiliary uniform-excess contraction bound]
\lean{Randomized.RegularMultigraph.uniformExcess_walkDist_two_mul_le_aux}
\label{Randomized.RegularMultigraph.uniformExcess_walkDist_two_mul_le_aux}
\leanok
\uses{Randomized.RegularMultigraph.Connected, Randomized.RegularMultigraph.IsBipartite, Randomized.RegularMultigraph.one_div_pow_le_transProb_two_mul, Randomized.RegularMultigraph.sum_transProb_left, Randomized.RegularMultigraph.transProb, Randomized.RegularMultigraph.walkDist, Randomized.uniformExcess, Randomized.uniformExcess_le_of_minorant}
\difficulty{4}
Let $G$ be a finite regular multigraph of positive degree $D$ with a nonempty vertex
set, and suppose that $G$ is connected and nonbipartite. Then $|V|/D^{2|V|}\le 1$, and
every rational-valued function $p$ on the vertices whose values sum to $1$ satisfies:
after $2|V|$ steps, its distribution has uniform excess at most
$\left(1-|V|/D^{2|V|}\right)$ times the uniform excess of $p$.
\end{theorem}

\begin{theorem}[Uniform excess contraction under random walks]
\lean{Randomized.RegularMultigraph.exists_uniformExcess_walkDist_le}
\label{Randomized.RegularMultigraph.exists_uniformExcess_walkDist_le}
\leanok
\uses{Randomized.RegularMultigraph.Connected, Randomized.RegularMultigraph.IsBipartite, Randomized.RegularMultigraph.uniformExcess_walkDist_mul_le, Randomized.RegularMultigraph.uniformExcess_walkDist_two_mul_le_aux, Randomized.RegularMultigraph.walkDist, Randomized.RegularMultigraph.walkDist_zero, Randomized.uniformExcess, Randomized.uniformExcess_nonneg}
\difficulty{5}
Let $G$ be a regular multigraph of degree $D>0$ on a finite, nonempty vertex set $V$.
Assume that $G$ is connected and not bipartite. Then there is a walk length $k$ such
that, for every function $p:V\to\mathbb{Q}$ with $\sum_{v\in V}p(v)=1$, the distribution
obtained by taking a uniformly random length-$k$ walk from a starting vertex sampled
according to $p$ satisfies
\[\max_{v\in V}\left(\sum_{u\in V}p(u)\Pr[\text{the walk from }u\text{ ends at }v]-\frac{1}{|V|}\right)\leq \left(1-\frac{1}{|V|^{10|V|}}\right)\max_{v\in V}\left(p(v)-\frac{1}{|V|}\right).\]
\end{theorem}

\begin{theorem}[Uniform-excess contraction on nonbipartite graphs]
\lean{Randomized.RegularMultigraph.uniformExcess_walkDist_two_mul_le}
\label{Randomized.RegularMultigraph.uniformExcess_walkDist_two_mul_le}
\leanok
\uses{Randomized.RegularMultigraph.Connected, Randomized.RegularMultigraph.IsBipartite, Randomized.RegularMultigraph.uniformExcess_walkDist_two_mul_le_aux, Randomized.RegularMultigraph.walkDist, Randomized.uniformExcess, Randomized.uniformExcess_nonneg}
\difficulty{4}
Let $G$ be a finite regular multigraph of positive degree $D$ with a nonempty vertex
set. Suppose $G$ is connected and nonbipartite, and $D\le |V|^5$, where $|V|$ is the
number of vertices. For every rational-valued function $p$ on the vertices whose values
sum to $1$, the distribution after $2|V|$ steps has uniform excess at most
$\left(1-1/|V|^{10|V|}\right)$ times the uniform excess of $p$.
\end{theorem}

\begin{theorem}[Stationarity under repeated walks]
\lean{Randomized.RegularMultigraph.IsStationary.walkDist_eq}
\label{Randomized.RegularMultigraph.IsStationary.walkDist_eq}
\leanok
\uses{Randomized.RegularMultigraph.IsStationary, Randomized.RegularMultigraph.walkDist, Randomized.RegularMultigraph.walkDist_succ, Randomized.RegularMultigraph.walkDist_zero}
\difficulty{4}
Let $G$ be a finite regular multigraph of positive degree, and let $p$ assign a rational
number to each vertex. If one step of the random walk preserves $p$, then for every
nonnegative integer $k$, the distribution obtained by applying the $k$-step walk
transition probabilities to $p$ is again $p$.
\end{theorem}

\begin{theorem}[Uniqueness of the stationary uniform distribution]
\lean{Randomized.RegularMultigraph.isStationary_iff_eq_uniform}
\label{Randomized.RegularMultigraph.isStationary_iff_eq_uniform}
\leanok
\uses{Randomized.RegularMultigraph.Connected, Randomized.RegularMultigraph.IsBipartite, Randomized.RegularMultigraph.IsStationary, Randomized.RegularMultigraph.IsStationary.walkDist_eq, Randomized.RegularMultigraph.exists_uniformExcess_walkDist_le, Randomized.RegularMultigraph.uniform_isStationary, Randomized.le_uniformExcess, Randomized.sum_sub_inv_card, Randomized.uniformExcess, Randomized.uniformExcess_nonneg}
\difficulty{5}
Let $G$ be a finite nonempty $D$-regular multigraph with $D>0$. Suppose $G$ is connected
and nonbipartite. For a rational-valued function $p$ on the vertices whose values sum to
$1$, call $p$ stationary if one step of the walk preserves it: for every vertex $v$,
$\sum_u p(u)P(u,v)=p(v)$, where $P(u,v)$ is the probability that a uniformly chosen
indexed neighbor of $u$ is $v$. Then $p$ is stationary if and only if $p(v)=1/|V(G)|$
for every vertex $v$.
\end{theorem}

\begin{theorem}[Real stationary distributions are uniform]
\lean{Randomized.RegularMultigraph.isStationary_real_iff_eq_uniform}
\label{Randomized.RegularMultigraph.isStationary_real_iff_eq_uniform}
\leanok
\uses{Randomized.RegularMultigraph.Connected, Randomized.RegularMultigraph.IsBipartite, Randomized.RegularMultigraph.one_div_pow_le_transProb_two_mul, Randomized.RegularMultigraph.sum_transProb_left, Randomized.RegularMultigraph.transProb, Randomized.RegularMultigraph.transProb_succ_eq_sum, Randomized.RegularMultigraph.transProb_zero}
\difficulty{6}
Let $G$ be a finite nonempty $D$-regular multigraph with $D>0$, and suppose it is
connected and nonbipartite. Let $p$ be a real-valued function on its vertices whose
values sum to $1$. For each pair of vertices $u,v$, let $P(u,v)$ be the probability that
a uniformly chosen indexed neighbor of $u$ is $v$. Then the condition $\sum_u
p(u)P(u,v)=p(v)$ for every vertex $v$ holds if and only if $p(v)=1/|V(G)|$ for every
vertex $v$.
\end{theorem}

\begin{theorem}[Asymptotic lower bound on expected visits]
\lean{Randomized.RegularMultigraph.exists_le_expVisits}
\label{Randomized.RegularMultigraph.exists_le_expVisits}
\leanok
\uses{Randomized.RegularMultigraph.Connected, Randomized.RegularMultigraph.IsBipartite, Randomized.RegularMultigraph.exists_uniformExcess_walkDist_le, Randomized.RegularMultigraph.expVisits, Randomized.RegularMultigraph.expVisits_eq, Randomized.RegularMultigraph.sum_walkDist, Randomized.RegularMultigraph.transProb, Randomized.RegularMultigraph.transProb_nonneg, Randomized.RegularMultigraph.uniformExcess_walkDist_le_self, Randomized.RegularMultigraph.uniformExcess_walkDist_mul_le, Randomized.RegularMultigraph.walkDist, Randomized.RegularMultigraph.walkDist_add, Randomized.inv_card_sub_le, Randomized.uniformExcess}
\difficulty{6}
Let $G$ be a regular multigraph on a finite, nonempty vertex set $V$, with positive
degree. Suppose every pair of vertices is connected by a finite sequence of neighbor
steps, and no Boolean coloring assigns opposite colors across every neighbor step. For
any vertices $u,v$, there is an integer $N_0$ such that for every integer $N\ge N_0$,
the expected number of visits to $v$ during the length-$N$ random walk started at $u$ is
at least $N/(2|V|)$.
\end{theorem}
